\documentclass[11pt,a4paper]{article}
\usepackage[T1]{fontenc}
\usepackage[utf8]{inputenc}
\usepackage{lmodern,amsmath,amssymb,amsthm,mathtools}
\usepackage[margin=25mm]{geometry}
\usepackage{microtype,booktabs,array}
\usepackage[hidelinks]{hyperref}
\usepackage{xurl}
\newtheorem{theorem}{Theorem}[section]
\newtheorem{lemma}[theorem]{Lemma}
\newtheorem{proposition}[theorem]{Proposition}
\newtheorem{corollary}[theorem]{Corollary}
\theoremstyle{definition}\newtheorem{hypothesis}[theorem]{Hypothesis}
\theoremstyle{remark}\newtheorem{remark}[theorem]{Remark}
\newcommand{\R}{\mathbb R}
\newcommand{\PP}{\mathbb P}
\newcommand{\Q}{\mathcal Q}
\newcommand{\G}{\mathcal G_3}
\newcommand{\dd}{\,\mathrm d}
\newcommand{\norm}[1]{\lVert #1\rVert}
\newcommand{\diver}{\operatorname{div}}
\numberwithin{equation}{section}
\allowdisplaybreaks[2]
\setlength{\parskip}{3pt}
\title{Dissipation-budget continuation\\and critical divergence-defect criteria}
\author{}
\date{6 September 2026}
\begin{document}
\maketitle
\begin{abstract}
Retaining cubic dissipation removes the endpoint continuation theorem from
the existing strict pressure-absorption route. We give its explicit enstrophy
bound, extend the quotient divergence-defect criterion from $L^4_tL^2_x$ to
every finite spatial exponent on the nonendpoint derivative-critical line,
and reconstruct the finite-horizon defect--enstrophy comparison in HF30.
The comparison is resolved for actual classical trajectories, not by a
pointwise reverse norm inequality. An additional endpoint-smallness statement
is proved with its strict threshold. None of these arguments supplies an
arbitrary-data critical estimate.
\end{abstract}
\noindent\fbox{\begin{minipage}{0.96\linewidth}
\textbf{Status.} Full component derivations, dependent on the explicitly stated
existing local and quotient packages. Author self-check only; independent
mathematical audit is pending. No claim-graph promotion and no claim of a
complete NS-R3 proof. This supplement is not a replacement for the manuscript.
\end{minipage}}

\section{Setting, provenance, and exact inputs}
We work only with the original unforced equation on $\R^3$,
\begin{equation}\label{db:NS}
 u_t+(u\cdot\nabla)u+\nabla p=\nu\Delta u,
 \qquad\diver u=0,\qquad u(0)=u_0,
\end{equation}
where $\nu>0$ and $u_0$ is real, solenoidal, and Schwartz. All norms are on
$\R^3$. The gradient convention is $(\nabla u)_{ij}=\partial_j u_i$.
For the selected maximal classical branch on $[0,T_*)$ set
\[
 E(t)=\norm{u(t)}_2^2,\quad Y(t)=\norm{\nabla u(t)}_2^2,
 \quad Z(t)=\norm{\Delta u(t)}_2^2,
 \qquad E_0=E(0),\quad Y_0=Y(0).
\]
The imported local package is precisely the manuscript's LOCAL node:
$u,p\in C^j([0,t];H^m)$ for every $j,m$ and $t<T_*$, uniqueness in the
$H^1$ mild class, and $Y(t)\to\infty$ at a finite $T_*$. The energy identity
\begin{equation}\label{db:energy}
 E(t)+2\nu\int_0^tY(s)\dd s=E_0
\end{equation}
is an existing manuscript result. This assembled package is a project input
based on Tao's local theory, not a theorem newly proved here \cite{paper,tao}.
No positive-time Schwartz persistence is assumed.

Fix the ordinary homogeneous Sobolev constant $S$ in
$\norm f_6\le S\norm{\nabla f}_2$. This also holds for finite arrays by
applying the scalar inequality to their magnitude. Riesz transforms and the
Leray projection have their usual unweighted $L^p$ bounds for $1<p<\infty$.
These are the manuscript's analytic foundations; no weighted multiplier
bound is assumed. Throughout,
\begin{equation}\label{db:constants}
 a_0=\frac98S^2,\qquad k=\frac{27}{16}S^2.
\end{equation}

\paragraph{Relation to earlier work.}
The nonendpoint mechanism is classical Serrin-type interpolation and an
enstrophy test \cite{serrin}; that test is proved below rather than imported
as a regularity theorem. HF28 already retained dissipation to avoid the
endpoint theorem in its comparison-field route. HF30 reconstructs the
fourth-power defect comparison \cite{hf28,hf30}. The present dependency
repair applies that principle directly to the manuscript's original strict
pressure hypothesis. The full finite-exponent defect calculation extends the
single exponent displayed in HF25/HF30. We make no novelty or priority claim.
Serrin's bibliographic metadata and Tao's primary preprint record were checked;
no unseen theorem from those sources is used to fill a missing estimate.

\paragraph{Frozen research inputs.}
The initial reads used \texttt{itpplasma/navier} at
\nolinkurl{9c5def29be04e0d0e723677feae0437a0ea61300} and
\texttt{itpplasma/navier-paper} at
\nolinkurl{065a962d6c6672dcf4bed2db17eb004081014a68}.
The companion evidence record distinguishes these read pins from the later
commit parents chosen after refreshing the remote branches.

\section{An elementary continuation lemma}
\begin{lemma}[Direct $L^4_tL^6_x$ continuation]\label{db:direct}
On every compact classical interval,
\begin{equation}\label{db:enstrophy}
 Y'+\nu Z\le k\nu^{-3}\norm u_6^4Y.
\end{equation}
For $I(t)=\nu^{-3}\int_0^t\norm{u(s)}_6^4\dd s$ this implies
\begin{equation}\label{db:YZ}
 Y(t)+\nu\int_0^t Z(s)\dd s\le Y_0e^{kI(t)}.
\end{equation}
If $H<\infty$ and $I(t)$ is bounded uniformly for
$0<t<\min(H,T_*)$, then $T_*>H$.
\end{lemma}
\begin{proof}
Test \eqref{db:NS} with $-\Delta u$. Pressure vanishes by solenoidality.
H\"older, interpolation, Sobolev, and Plancherel give
\[
 \frac12Y'+\nu Z
 \le\norm u_6\norm{\nabla u}_3\norm{\Delta u}_2
 \le S^{1/2}\norm u_6Y^{1/4}Z^{3/4}.
\]
For $b,x\ge0$ and $\epsilon>0$, maximization in $x$ gives
\[
 bx^{3/4}\le\epsilon x+\frac{27b^4}{256\epsilon^3}.
\]
Use $\epsilon=\nu/2$ and multiply by two. This proves
\eqref{db:enstrophy}, including its constant. An integrating factor yields
\[
 Y(t)+\nu\int_0^t e^{k(I(t)-I(s))}Z(s)\dd s\le Y_0e^{kI(t)}.
\]
Since $I$ is nondecreasing, \eqref{db:YZ} follows. All spatial tests are first
justified on compact classical intervals; the all-orders Sobolev package
makes the boundary products integrable and cutoff errors vanish. Uniformity
of the displayed bound, together with \eqref{db:energy}, contradicts the
local blow-up alternative if $T_*\le H$. This includes $T_*=H$; no solution
at the hypothetical endpoint has been presumed.
\end{proof}

\section{Strict pressure absorption without the endpoint theorem}
Let $r=|u|$ and use the weak gradient of $r$, with zero-set conventions as in
the manuscript. Define
\begin{equation}\label{db:D3}
 X=\frac13\norm u_3^3,\qquad
 D_3=\int_{\R^3}\bigl(r|\nabla u|^2+r|\nabla r|^2\bigr)\dd x,
 \qquad P_3=\int_{\R^3}p\,u\cdot\nabla r\dd x.
\end{equation}
The imported integrated cubic balance is
\begin{equation}\label{db:cubic}
 X(t)+\nu\int_0^tD_3(s)\dd s=X(0)+\int_0^tP_3(s)\dd s.
\end{equation}

\begin{lemma}[Cubic dissipation controls the nonendpoint norm]\label{db:speed}
At every classical time,
\begin{equation}\label{db:l6speed}
 \norm u_9^3\le a_0D_3,\qquad
 \norm u_6^4\le a_0\norm u_3D_3.
\end{equation}
\end{lemma}
\begin{proof}
The Lipschitz chain rule gives $|\nabla r|\le|\nabla u|$ almost everywhere.
On a compact classical interval, $u$ is bounded and in $H^1\cap L^3$;
therefore $r^{3/2}\in H^1$ and
$\nabla(r^{3/2})=(3/2)r^{1/2}\nabla r$ almost everywhere. At zeros the
factor $r^{1/2}$ makes this formula unambiguous. Consequently
\[
 D_3\ge2\int r|\nabla r|^2
       =\frac89\norm{\nabla(r^{3/2})}_2^2.
\]
Sobolev applied to $r^{3/2}$ gives the first inequality. Interpolating
$\norm u_6\le\norm u_3^{1/4}\norm u_9^{3/4}$ gives the second. This
argument introduces no derivative divided by $|u|$ at its zero set.
\end{proof}

\begin{theorem}[Quantitative direct pressure-absorption suffix]\label{db:pressure}
Fix a finite $H>0$. Suppose $0\le\theta<1$ and $A\ge0$ are finite and
\begin{equation}\label{db:absorption}
 \int_0^t P_3(s)\dd s\le\theta\nu\int_0^tD_3(s)\dd s+A
 \quad\text{for every }0<t<\min(H,T_*).
\end{equation}
Put $\delta=1-\theta$ and $R=\norm{u_0}_3^3+3A$. Then
\begin{align}
 \norm{u(t)}_3^3+3\delta\nu\int_0^tD_3\dd s&\le R,\label{db:R}\\
 \int_0^t\norm u_6^4\dd s&\le\frac{3S^2}{8\delta\nu}R^{4/3},
 \label{db:pressurel6}\\
 Y(t)+\nu\int_0^tZ\dd s
 &\le Y_0\exp\left(\frac{81S^4R^{4/3}}{128\delta\nu^4}\right).
 \label{db:pressureY}
\end{align}
In particular $T_*>H$. If the hypothesis is supplied for every input datum,
viscosity, and finite horizon, it implies the programme's global NS-R3
conclusion using LOCAL and ENERGY, without ESS.
\end{theorem}
\begin{proof}
Subtract \eqref{db:absorption} from \eqref{db:cubic} to obtain
\eqref{db:R}. Thus $\sup\norm u_3\le R^{1/3}$ and
$\int D_3\le R/(3\delta\nu)$. Integrating \eqref{db:l6speed} proves
\eqref{db:pressurel6}, since $a_0/3=3S^2/8$. Lemma~\ref{db:direct}
then proves \eqref{db:pressureY} and $T_*>H$. A finite maximal time would
contradict the conclusion for a horizon larger than it. Global smoothness and
the uniform kinetic-energy bound follow from the imported local package and
\eqref{db:energy}. The datum $u_0=0$ is handled by uniqueness; no division by
$Y_0$ or $R$ has been used.
\end{proof}

\begin{corollary}[The original high-pressure hypothesis suffices directly]
In the existing split $P_3=Q_J+L_J$, suppose the manuscript's low-frequency
lemma supplies
$\int_0^tL_J\dd s\le L_{J,H}<\infty$ uniformly below $\min(H,T_*)$,
and HIGH-PRESSURE supplies
$\int_0^tQ_J\dd s\le\theta\nu\int_0^tD_3\dd s+A_{J,H}$
with $\theta<1$. Theorem~\ref{db:pressure} applies with
$A=A_{J,H}+L_{J,H}$. The cutoff and remainders retain exactly their existing
input-dependent quantifiers; no new smallness condition is imposed.
\end{corollary}

\begin{remark}[What cannot be deleted]
This replaces the \emph{sufficiency path from strict absorption}, not the
standalone theorem that a mere $L^\infty_tL^3_x$ bound implies continuation.
The latter does not supply $\int D_3$ and retains its existing endpoint input.
Likewise a high-strain hypothesis allowing $\theta=1$ must not be silently
strengthened to $\theta<1$. At the scalar level, $X(t)=1$,
$D(t)=(H-t)^{-1}$ and $K(t)=\nu D(t)$ obey $X'+\nu D=K$ and the
$\theta=1$ bound with zero remainder, while $\int_0^H D=\infty$.
These are not Navier--Stokes fields; they isolate the invalid inference from
coefficient-one absorption to a dissipation budget.
\end{remark}

\section{The quotient package and its exact mixed-pressure identity}
Set
\[
 \G=\overline{\{\nabla\phi:\phi\in C_c^\infty(\R^3)\}}^{L^3},\qquad
 w=\arg\min_{z\in u+\G}\frac13\norm z_3^3,
 \quad q=w-u,\quad Q=\Q(u)=\frac13\norm w_3^3.
\]
Write $A_w=|w|w$, $\sigma=-\diver w=-\diver q$, and let $D$ denote the
quotient dissipation. These are not the physical pressure or the velocity's
ordinary divergence. The following precise statements are imported from the
manuscript's audited quotient and unweighted div--curl package
\cite{paper,hf25}:
\begin{align}
 &u=\PP w,\quad\diver A_w=0,\quad
 Q'+\nu D=-\int q_i u_j\partial_j(A_w)_i,\label{db:Qidentity}\\
 &\norm w_9^3\le a_0D,\quad A_w\in W^{1,3/2},\quad
 w,q\in L^6\cap W^{1,2}_{\mathrm{loc}},\label{db:Qtoolkit}\\
 &\norm{\nabla q}_2^2=\norm\sigma_2^2\le\frac14Y.
 \label{db:defectspatial}
\end{align}
$Q$ is locally absolutely continuous in time. The derivative formulas use
global weak derivatives across zero sets; $w,q\in L^2$ is not asserted.
These are project results, not relabelled published literature theorems.

Let $R_i$ have Fourier symbol $-i\xi_i/|\xi|$. For $1<s<\infty$ define a
finite unweighted multiplier constant $M_s$ by
\begin{equation}\label{db:Ms}
 \left\|\bigl(R_iR_j f\bigr)_{i,j=1}^3\right\|_s\le M_s\norm f_s.
\end{equation}
Plancherel permits $M_2=1$. Write $C_p=\norm{\PP}_{L^p\to L^p}$.

\begin{lemma}[Mixed identity with no extra defect assumption]\label{db:mixed}
At classical times,
\begin{equation}\label{db:mixedformula}
 K:=Q'+\nu D=\int (R_iR_j\sigma)(A_w)_i u_j.
\end{equation}
In particular the identity is already an $L^2$ pairing before any additional
$L^s$ hypothesis on $\sigma$ is imposed.
\end{lemma}
\begin{proof}
Integrating \eqref{db:Qidentity} by parts and using $\diver u=0$ gives
$K=\int (A_w)_i u_j\partial_jq_i$. The weak gradient of $q$ is symmetric,
so $\partial_jq_i=\partial_iq_j$. Since $q$ is a distributional gradient
and $\sigma=-\diver q$, the Fourier identity is
$\partial_iq_j=R_iR_j\sigma$. A possible distribution supported at zero
frequency disappears because both derivative fields are in $L^2$.
The products are integrable: $uA_w\in L^2$ by $u,w\in L^6$, while
$\nabla q\in L^2$. In the integration-by-parts step, the cutoff boundary
product is integrable by boundedness of $u$, $q\in L^3$, and
$A_w\in L^{3/2}$. The original differentiated product is integrable by
\eqref{db:Qtoolkit} and the existing quotient identity. Density and removal
of cutoffs therefore prove \eqref{db:mixedformula}.
\end{proof}

\section{The entire finite-exponent nonendpoint defect line}
\begin{theorem}[Critical scalar-defect criterion]\label{db:family}
Let $3/2<s<\infty$ and define
\begin{equation}\label{db:exponents}
 \alpha=\frac{3}{2s},\quad p_s=\frac{2s}{2s-3}=\frac1{1-\alpha},
 \quad r_s=\frac{3s}{s-1},\quad
 b_s=M_s C_{r_s}3^{1-\alpha}a_0^\alpha,
\end{equation}
\begin{equation}\label{db:cs}
 c_s=(1-\alpha)\alpha^{\alpha/(1-\alpha)}
           2^{\alpha/(1-\alpha)}b_s^{p_s},\qquad
 h_s(t)=c_s\nu^{-3/(2s-3)}\norm{\sigma(t)}_s^{p_s}.
\end{equation}
Suppose, for a finite horizon $H$, that $\sigma(t)\in L^s$ almost everywhere
and
\begin{equation}\label{db:Hfinite}
 \sup_{t<\min(H,T_*)}\mathcal H_s(t)<\infty,\qquad
 \mathcal H_s(t)=\int_0^t h_s(\tau)\dd\tau.
\end{equation}
Then $T_*>H$. Quantitatively, with
\begin{equation}\label{db:J0}
 B_6=3^{1/3}C_6^4a_0,\qquad
 J_0=\frac32 B_6\frac{Q(0)^{4/3}}{\nu^4},
\end{equation}
we have
\begin{align}
 Q(t)&\le Q(0)e^{\mathcal H_s(t)},\label{db:Qbound}\\
 \nu^{-3}\int_0^t\norm u_6^4\dd\tau
 &\le J_0e^{4\mathcal H_s(t)/3},\label{db:familyl6}\\
 Y(t)+\nu\int_0^t Z\dd\tau
 &\le Y_0\exp\bigl(kJ_0e^{4\mathcal H_s(t)/3}\bigr).
 \label{db:familyY}
\end{align}
\end{theorem}
\begin{proof}
First apply Lemma~\ref{db:mixed}, H\"older, and the Leray bound. Since
$3<r_s<9$, all norms of $w$ below are supplied by interpolation between the
existing $L^3$ and $L^9$ bounds:
\begin{align*}
 |K|&\le M_s\norm\sigma_s\norm u_{r_s}\norm w_{r_s}^2\\
 &\le M_s C_{r_s}\norm\sigma_s\norm w_{r_s}^3
 \le b_s\norm\sigma_s Q^{1-\alpha}D^\alpha.
\end{align*}
The interpolation exponent is exactly $\alpha$, because
$1/r_s=(1-\alpha)/3+\alpha/9$. For $0<\alpha<1$, maximizing
$bx^\alpha-\epsilon x$ gives
\begin{equation}\label{db:younggeneral}
 bx^\alpha\le\epsilon x+
 (1-\alpha)\alpha^{\alpha/(1-\alpha)}
 \epsilon^{-\alpha/(1-\alpha)}b^{1/(1-\alpha)}.
\end{equation}
Use $b=b_s\norm\sigma_s Q^{1-\alpha}$ and $\epsilon=\nu/2$ to obtain
\begin{equation}\label{db:defectODE}
 Q'+\frac\nu2D\le h_s Q.
\end{equation}
The $L^2$ and $L^s$ Riesz extensions agree on $L^2\cap L^s$ by smooth
approximation. Thus the $s=2$ pairing and the $s$-pairing represent the same
integrable function; no new integration-by-parts step in an unproved
regularity class is needed.
The coefficient $h_s$ is measurable, for instance using the countable-test
characterization of the $L^s$ norm of the distribution $\sigma$. The norm
of an absent $L^s$ representative is understood as infinity.

Gronwall gives \eqref{db:Qbound}. For the sharper dissipation estimate,
multiply \eqref{db:defectODE} by $(4/3)Q^{1/3}$:
\[
 (Q^{4/3})'+\frac{2\nu}{3}Q^{1/3}D
 \le\frac43 h_sQ^{4/3}.
\]
This is valid at $Q=0$ because the power is continuously differentiable there;
there is no division by $Q$. Its integrating factor gives
\[
 Q(t)^{4/3}+\frac{2\nu}{3}\int_0^t
 e^{\frac43(\mathcal H_s(t)-\mathcal H_s(\tau))}
 Q(\tau)^{1/3}D(\tau)\dd\tau
 \le Q(0)^{4/3}e^{4\mathcal H_s(t)/3}.
\]
The exponential in the integral is at least one. Moreover
\[
 \norm u_6^4\le C_6^4\norm w_3\norm w_9^3
              \le B_6Q^{1/3}D.
\]
These two inequalities prove \eqref{db:familyl6}, with exactly the constant
in \eqref{db:J0}. Lemma~\ref{db:direct} proves \eqref{db:familyY} and the
continuation conclusion. Every step is made first below the maximal time;
uniformity in \eqref{db:Hfinite} is the input excluding its endpoint.
\end{proof}

\begin{remark}[Critical scaling, the endpoint, and the missing bound]
Under Navier--Stokes dilation $u_\lambda(t,x)=\lambda u(\lambda^2t,\lambda x)$,
uniqueness of the minimizer gives
$w_\lambda=\lambda w(\lambda^2t,\lambda x)$ and
$\sigma_\lambda=\lambda^2\sigma(\lambda^2t,\lambda x)$.
Thus the time-space integral in \eqref{db:Hfinite} is invariant precisely
when $2/p_s+3/s=2$, as satisfied here. The theorem does not cover $s=\infty$;
no $L^1$ or $L^\infty$ endpoint Riesz bound was used. The case $s=3/2$
requires the distinct smallness statement below. Most importantly,
\eqref{db:Hfinite} is a hypothesis, not an estimate established for arbitrary
data. Changing exponents has not solved the producer problem.
\end{remark}

\begin{proposition}[Endpoint defect smallness]\label{db:endpoint}
Put $b_*=M_{3/2}C_9a_0$. If, uniformly below $\min(H,T_*)$,
\[
 b_*\norm\sigma_{3/2}\le\theta\nu\quad\text{almost everywhere},
 \qquad 0\le\theta<1,
\]
then $Q'+(1-\theta)\nu D\le0$ and
\[
 \int_0^t\norm u_6^4\dd\tau
 \le\frac{B_6Q(0)^{4/3}}{(1-\theta)\nu}.
\]
In particular $T_*>H$.
\end{proposition}
\begin{proof}
Lemma~\ref{db:mixed} with $s=3/2$ gives
$|K|\le M_{3/2}C_9\norm\sigma_{3/2}\norm w_9^3
\le b_*\norm\sigma_{3/2}D$.
Hence $Q(t)+(1-\theta)\nu\int_0^tD\le Q(0)$. Multiply the bound on
$\int D$ by $B_6Q(0)^{1/3}$ and use Lemma~\ref{db:direct}.
No assertion is made for a merely finite, arbitrarily large
$L^\infty_tL^{3/2}_x$ defect norm.
\end{proof}

\section{Resolution of the finite-horizon comparison in HF30}
Define
\[
 \Gamma(t)=\nu^{-3}\int_0^t\norm\sigma_2^4\dd\tau,
 \qquad B(t)=\nu^{-3}\int_0^tY(\tau)^2\dd\tau,
 \qquad C_\sigma=\frac{81}{32}C_6^4a_0^3.
\]
\begin{corollary}[Defect and enstrophy integrability on the actual branch]
\label{db:comparison}
For $t<T_*$,
\begin{equation}\label{db:comparisonbounds}
 \Gamma(t)\le\frac1{16}B(t),\qquad
 B(t)\le\frac{E_0Y_0}{2\nu^4}
   \exp\bigl(kJ_0e^{4C_\sigma\Gamma(t)/3}\bigr).
\end{equation}
For any $0<L\le T_*$ with $L<\infty$, finiteness of
$\int_0^L\norm\sigma_2^4$ is equivalent to finiteness of
$\int_0^L Y^2$. If $L=T_*<\infty$, both integrals must be infinite.
\end{corollary}
\begin{proof}
At $s=2$, \eqref{db:exponents} gives $\alpha=3/4$, $p_s=4$, and $r_s=6$.
With $M_2=1$, the constant in \eqref{db:cs} is $C_\sigma$; thus
$\mathcal H_2=C_\sigma\Gamma$ and Theorem~\ref{db:family} applies on every
compact classical interval. The first inequality of
\eqref{db:comparisonbounds} follows by squaring
\eqref{db:defectspatial}. The second follows from
\[
 B(t)\le\nu^{-3}\sup_{0\le\tau\le t}Y(\tau)
                \int_0^tY(\tau)\dd\tau
       \le\frac{E_0}{2\nu^4}\sup_{0\le\tau\le t}Y(\tau)
\]
and \eqref{db:familyY}. Nondecreasing $\Gamma$ makes the bound at time $t$
a common bound for all earlier times. Monotone convergence proves the
finiteness equivalence at $L$. Finiteness at $T_*<\infty$ contradicts
Lemma~\ref{db:direct}. This reconstructs the HF30 comparison using the
already stated project inputs; it is not a new unconditional regularity
assertion.
\end{proof}

\paragraph{The precise logical repair.}
The current DEFECT-L4 review says the finite-lifespan comparison is
``unsettled in both directions.'' Corollary~\ref{db:comparison} resolves it
under the same actual-branch hypotheses. This is not a reverse instantaneous
inequality $Y^2\lesssim\norm\sigma_2^4$, which has not been proved. Nor does
$\norm\sigma_2\le\norm{\nabla u}_2/2$ by itself show that the defect
criterion is not weaker: it says that the gradient hypothesis implies the
defect hypothesis. The converse in this corollary uses evolution and
continuation, not that one-sided spatial bound. There is no strictly weaker
finite-endpoint trajectory class exhibited by this comparison.

\paragraph{Why energy still does not close it.}
The available unconditional estimate is only
\[
 \int_0^{T_*}\norm\sigma_2^2\dd t\le\frac{E_0}{8\nu}.
\]
For an elementary integrability check, $g(t)=(L-t)^{-1/3}$ has finite
$\int_0^Lg^2$ but infinite $\int_0^Lg^4$. It is not an actual defect or a
Navier--Stokes solution. It simply prevents replacing the required fourth
power by the energy-controlled second power without an additional argument.

\section{Dependency repair and remaining proof task}
For the strict pressure route the candidate dependency path is
\[
 \begin{gathered}
 \text{LOCAL, ENERGY, PRESSURE, LOW-PRESSURE, HIGH-PRESSURE}\\
 \Longrightarrow\ \text{Theorem~\ref{db:pressure}}
 \Longrightarrow\ \text{NS-R3}.
 \end{gathered}
\]
The nonendpoint Sobolev and enstrophy steps are proved in this supplement;
ESS, the Leray--Hopf extension, and the endpoint $L^3$ continuation theorem
are not on this sufficiency path. They remain dependencies of the separately
stated endpoint result. Thus this offers an alternative formalization target
without falsely declaring the entire existing manuscript independent of ESS.

For the defect route the path is the existing quotient package, the
mixed identity, \eqref{db:defectODE}, the weighted $Q^{4/3}$ balance, and
Lemma~\ref{db:direct}. Again there is no endpoint theorem. The missing
arbitrary-data input is an upper bound on at least one critical accumulated
defect norm, for instance $\Gamma$, uniformly to a putative finite maximal
time. Neither the exponent family nor the comparison provides that bound.
The spectral-index producer of HF30 is unchanged.

The next integration steps are a separate mathematical audit of these
component derivations, the scoped repair of DEFECT-L4's review text, and an
additional conditional graph edge for strict absorption. Do not delete ESS
from the bare-CRITICAL theorem, do not promote any gap node, and do not treat
improved constants or a conditional equivalence as an arbitrary-data estimate.

\begin{thebibliography}{9}
\bibitem{paper} \texttt{itpplasma/navier-paper}, \emph{main.tex}, revision
\nolinkurl{065a962d6c6672dcf4bed2db17eb004081014a68}.
Project-level input: local package, integrated cubic balance, quotient
functional and div--curl results. Private source, not published literature.
\bibitem{hf25} \texttt{itpplasma/navier},
\nolinkurl{research/evidence/hf25-beyond-hf21-continuation.md} and
\nolinkurl{hf25-review-defect-criterion.md}, as present at the research read pin
in Section~1. Project evidence and recorded audit, not an external theorem.
\bibitem{hf28} \texttt{itpplasma/navier},
\nolinkurl{research/evidence/hf28-weighted-spectral-continuation.tex},
as present at the research read pin in Section~1. Prior project continuation.
\bibitem{hf30} \texttt{itpplasma/navier},
\nolinkurl{research/evidence/hf30-second-order-continuation.tex},
as present at the research read pin in Section~1, especially its section ``A quantitative resolution of the defect--enstrophy comparison.''
Imported candidate; no independent audit claimed by this supplement.
\bibitem{tao} T.~Tao, \emph{Localisation and compactness properties of the
Navier--Stokes global regularity problem}, primary preprint arXiv:1108.1165.
\url{https://arxiv.org/abs/1108.1165}.
The local/maximal package is imported from the manuscript's assembled proof,
not inferred from the preprint abstract.
\bibitem{serrin} J.~Serrin, \emph{On the interior regularity of weak solutions
of the Navier--Stokes equations}, Archive for Rational Mechanics and Analysis
\textbf{9} (1962), 187--195. \url{https://doi.org/10.1007/BF00253344}.
Historical attribution and verified metadata only; the nonendpoint estimate
used here is proved in Lemma~\ref{db:direct}.
\end{thebibliography}
\end{document}
